% !TEX root = ../main.tex
\chapter{ケーススタディ：List 全体のマイグレーションと関手則}
\label{chap:ch30-list-migration-functor-laws}
\BookIndex{ばっちいこう}{バッチ移行}
\BookIndex{List.map}{\texttt{List.map}}
\BookIndex{かんしゅそく}{関手則}

\begin{chapterabstract}
本章では、1 件のデータ移行を \texttt{List} 全体の移行へ持ち上げる方法を扱います。
実務では、\texttt{UserV1} から \texttt{UserV2} へ 1 件だけ変換できても、それだけでは十分ではありません。
通常は、ユーザー一覧、注文一覧、ログ一覧のような複数件のデータを一括で移行します。
このとき自然に現れるのが、\texttt{List.map migrate} です。

圏論の言葉では、\texttt{List} は「型と関数の世界」における関手として読めます。
関手は、恒等関数と合成を保存します。
ソフトウェアの言葉に直すと、\texttt{map} はリストの形を壊さず、各要素だけを同じ規則で変換します。
本章では、件数保存、合成保存、round-trip の持ち上げを Lean 4 で確認します。
\end{chapterabstract}

\begin{learninggoals}
\item 1 件の移行関数を \texttt{List.map} によってバッチ移行へ持ち上げる考え方を説明できます。
\item \texttt{List.map} がリストの長さを保存することを Lean で証明できます。
\item \texttt{List.map} の合成保存を、関手則として読めます。
\item 1 件の round-trip 性質を、リスト全体の round-trip 性質へ拡張できます。
\item バッチ移行で証明できる性質と、実システム側で別途確認すべき性質を区別できます。
\end{learninggoals}

\section{1件の保証をコレクション全体へ広げる}

\ChapterRef{chap:ch28_user_schema_lossless_migration}では、\texttt{UserV1} と \texttt{UserV2} の間の lossless migration を、1 件のデータに対して扱いました。
そこでは、\texttt{migrate : UserV1 -> UserV2} と \texttt{rollback : UserV2 -> UserV1} を定義し、変換して戻すと元に戻ることを示しました。
しかし、実務の移行対象は、たいてい 1 件ではありません。
DB のテーブル、CSV、イベントログ、API レスポンスの配列など、移行はほとんどの場合、コレクション全体に対して行われます。

ここで確認したいことは、単に「全件に対して変換関数を呼ぶ」ことではありません。
移行後も件数が変わらないでしょうか。
要素の順序は保たれるでしょうか。
1 件ごとの round-trip は、全件に対しても成り立つでしょうか。
二段階の移行をまとめてよいでしょうか。
これらは、バッチ移行のレビューでよく出てくる問いです。

図\ref{fig:ch30-overview}は、1 件の変換 \texttt{migrate} を各要素へ適用し、要素数と順序を保ったリストを作る流れを示します。

\FigurePlaceholder{fig-ch30-overview.png}{0.90}{移行のリストへの持ち上げ}{fig:ch30-overview}

\section{1件の移行から N 件の移行へ}

まず、1 件の移行を思い出します。
\texttt{UserV1} から \texttt{UserV2} への移行関数を \texttt{migrate} とします。
この関数は、1 人のユーザーを新しい形式へ変換します。

\[
  \texttt{migrate : UserV1 -> UserV2}
\]

バッチ移行では、入力は 1 件ではなく \texttt{List UserV1} になります。
出力は \texttt{List UserV2} です。
自然な定義は次のとおりです。

\[
  \texttt{migrateAll xs = xs.map migrate}
\]

これは、\texttt{xs} の各要素に \texttt{migrate} を適用し、その結果を同じ順序で並べたリストを返します。
要素を増やしたり減らしたりしません。
フィルタもしません。
グループ化もしません。
本章で扱うのは、このような「要素ごとの純粋な変換」です。

\texttt{List.map migrate} は純粋な変換核のモデルであり、ファイル入出力、トランザクション、失敗時スキップなどは含みません。

\section{List.map と関手の見方}

\texttt{List.map} は、リストの中身だけを変換します。
型 \texttt{A} から型 \texttt{B} への関数 \texttt{f : A -> B} があるとき、\texttt{List.map f} は \texttt{List A} から \texttt{List B} への関数になります。

\[
  \texttt{f : A -> B}
  \quad\Longrightarrow\quad
  \texttt{List.map f : List A -> List B}
\]

圏論の言葉では、これは \texttt{List} が関手として振る舞っているということです。
関手は、対象だけでなく射も移します。
ここでは、型 \texttt{A} を \texttt{List A} へ移し、関数 \texttt{f} を \texttt{List.map f} へ移しています。

本章で使う範囲では、関手則は次の二つとして読めます。
\begin{itemize}
\item 恒等保存：何もしない関数を \texttt{map} しても、リスト全体として何も変わりません。
\item 合成保存：\texttt{f} の後に \texttt{g} を適用する変換を一度に \texttt{map} しても、先に \texttt{f} を \texttt{map} し、次に \texttt{g} を \texttt{map} しても、結果は同じです。
\end{itemize}

数式風に書けば、合成保存は次の形です。

\[
  \texttt{map (g . f) = map g . map f}
\]

ここで \texttt{g . f} は、\texttt{f} を実行した後に \texttt{g} を実行する合成です。
この法則があるため、1 件の移行を組み合わせる設計と、リスト全体の移行を組み合わせる設計を対応させられます。

図\ref{fig:ch30-map-as-functor}は、型をリスト型へ、関数を \texttt{List.map} した関数へ対応させます。
この対応が恒等射と合成を保存することを、後で証明します。

\FigurePlaceholder{fig-ch30-map-as-functor.png}{0.90}{\texttt{List} 関手}{fig:ch30-map-as-functor}

\section{Lean 4 でドメインモデルを書く}

ここから、Lean 4 の小さなモデルを書きます。
例として、古いユーザースキーマ \texttt{UserV1} と新しいユーザースキーマ \texttt{UserV2} を用意します。
この例では、フィールド名とまとめ方が変わるだけで、情報は失われないものとします。

\begin{listing}[htbp]
\begin{minted}{lean4}
structure UserV1 where
  id : Nat
  name : String
  active : Bool
deriving Repr, DecidableEq

structure UserV2 where
  userId : Nat
  profileName : String
  enabled : Bool
deriving Repr, DecidableEq

def migrate (u : UserV1) : UserV2 :=
  { userId := u.id,
    profileName := u.name,
    enabled := u.active }

def rollback (v : UserV2) : UserV1 :=
  { id := v.userId,
    name := v.profileName,
    active := v.enabled }

def migrateAll (xs : List UserV1) : List UserV2 :=
  xs.map migrate

def rollbackAll (xs : List UserV2) : List UserV1 :=
  xs.map rollback
\end{minted}
\caption{\texttt{migrateAll} と \texttt{rollbackAll}}
\label{lst:ch30-domain-model}
\end{listing}

\texttt{migrateAll} は、新しい再帰関数として手で書くこともできます。
しかし、本章では意図的に \texttt{List.map} を使います。
\texttt{List.map} がすでに「各要素に同じ関数を適用し、リストの形を保つ」という意図を持っているためです。
実装上も仕様上も、この意図を読み取りやすくなります。

\section{件数保存}

バッチ移行で最初に確認したい性質は、件数保存です。
\texttt{migrateAll} は要素を変換するだけなので、入力リストと出力リストの長さは同じでなければなりません。

\[
  \texttt{(migrateAll xs).length = xs.length}
\]

この性質は、\texttt{List.map} の基本的な性質です。
Lean では、リストに対する帰納法で証明できます。
空リストなら、明らかに長さは 0 のままです。
先頭要素を一つ追加したリストでは、先頭を変換しても要素数は一つ増えたままであり、残りの部分については帰納法の仮定を使います。

\begin{listing}[htbp]
\begin{minted}{lean4}
theorem rollback_migrate (u : UserV1) :
    rollback (migrate u) = u := by
  cases u
  rfl

theorem migrate_rollback (v : UserV2) :
    migrate (rollback v) = v := by
  cases v
  rfl

theorem migrateAll_length (xs : List UserV1) :
    (migrateAll xs).length = xs.length := by
  induction xs with
  | nil =>
      rfl
  | cons x xs ih =>
      simp [migrateAll]
\end{minted}
\caption{\texttt{migrateAll\_length}}
\label{lst:ch30-length-preservation}
\end{listing}

\texttt{migrateAll\_length} は、バッチ移行に対する最小限の保証です。
ただし、件数が同じであることは、移行が正しいことの十分条件ではありません。
全件が壊れた値に置き換わっていても、件数だけは一致する可能性があります。
そのため、件数保存は入口であり、ID 保存や round-trip などの性質と組み合わせて使う必要があります。

\section{関手則1：恒等保存}

恒等保存は、\texttt{map} が余計な変化を加えないことを表します。
任意のリスト \texttt{xs} に対して、各要素に「何もしない関数」を適用しても、結果は元のリストと同じです。

\[
  \texttt{xs.map id = xs}
\]

ソフトウェアの言葉では、これは「no-op な一括変換は、バッチ全体としても no-op である」という主張です。
当然に見えますが、関手則として明示しておくと、adapter や移行ステップを足し引きするときの基準になります。

\section{関手則2：合成保存}

合成保存は、バッチ移行を設計するときに特に重要です。
たとえば、\texttt{UserV1 -> UserV2} の移行の後に、\texttt{UserV2 -> UserV3} の正規化を行うとします。
このとき、次の二つの方法が考えられます。

\begin{itemize}
\item 1 件ごとに「移行してから正規化する」関数を作り、それをリスト全体に \texttt{map} します。
\item 先にリスト全体を \texttt{UserV2} へ移行し、その後でリスト全体に正規化を \texttt{map} します。
\end{itemize}

関手の合成保存は、この二つが同じ結果になることを保証します。
Lean では、次のような一般形を証明できます。

\begin{listing}[htbp]
\begin{minted}{lean4}
theorem list_map_id {A : Type} (xs : List A) :
    xs.map (fun x => x) = xs := by
  induction xs with
  | nil =>
      rfl
  | cons x xs ih =>
      simp [ih]

theorem list_map_comp {A B C : Type}
    (f : A -> B) (g : B -> C) (xs : List A) :
    xs.map (fun x => g (f x)) = (xs.map f).map g := by
  induction xs with
  | nil =>
      rfl
  | cons x xs ih =>
      simp [ih]
\end{minted}
\caption{\texttt{list\_map\_id} と \texttt{list\_map\_comp}}
\label{lst:ch30-list-functor-laws}
\end{listing}

この証明も、リストに対する帰納法です。
空リストでは、どちらの方法でも空リストになります。
先頭要素を持つリストでは、先頭に対する変換は両辺で同じになり、残りのリストについては帰納法の仮定を使います。

合成保存は、\texttt{map} の中の関数が純粋な値変換としてモデル化されているときの性質です。
実装側でログ出力、DB 更新、外部 API 呼び出し、ランダム値生成などの副作用を含む場合、単純な \texttt{List.map} の法則だけでは実務上の振る舞いを説明できません。
その場合は、\texttt{Option}、\texttt{Except}、状態、ログなどを含む別のモデルが必要になります。

\section{二段階移行の合成保存}

一般形だけでは抽象的に見えるので、ユーザー移行の例に戻します。
\texttt{UserV2} から \texttt{UserV3} への正規化を考えます。
ここでは、\texttt{enabled} をそのまま \texttt{canLogin} として移す単純な例にします。

\begin{listing}[htbp]
\begin{minted}{lean4}
structure UserV3 where
  accountId : Nat
  displayName : String
  canLogin : Bool
deriving Repr, DecidableEq

def normalize (v : UserV2) : UserV3 :=
  { accountId := v.userId,
    displayName := v.profileName,
    canLogin := v.enabled }

def migrateThenNormalize (u : UserV1) : UserV3 :=
  normalize (migrate u)

def migrateAllThenNormalize (xs : List UserV1) : List UserV3 :=
  xs.map migrateThenNormalize

def normalizeAfterMigrateAll (xs : List UserV1) : List UserV3 :=
  (migrateAll xs).map normalize

theorem migrateAll_composition (xs : List UserV1) :
    migrateAllThenNormalize xs = normalizeAfterMigrateAll xs := by
  induction xs with
  | nil => rfl
  | cons x xs ih =>
      simp [migrateAllThenNormalize, normalizeAfterMigrateAll,
        migrateThenNormalize, migrateAll]
\end{minted}
\caption{\texttt{migrateAll\_composition}}
\label{lst:ch30-two-step-migration}
\end{listing}

この定理は、二段階の移行を「1 件単位で合成してからバッチ化する」方法と、「バッチ化したあと次の変換をかける」方法が一致することを述べています。
実務では、処理ステージをどこでまとめてもデータ変換の意味が変わらないというレビュー根拠になります。

図\ref{fig:ch30-composition-law}では、要素関数を先に合成する経路と、二回 \texttt{map} する経路が同じリストに到達します。

\FigurePlaceholder{fig-ch30-composition-law.png}{0.90}{合成保存}{fig:ch30-composition-law}

\section{round-trip をリスト全体へ持ち上げる}

\ChapterRef{chap:ch28_user_schema_lossless_migration}で扱った lossless migration では、1 件について次の二つを証明しました。

\[
  \texttt{rollback (migrate u) = u}
\]

\[
  \texttt{migrate (rollback v) = v}
\]

では、リスト全体ではどうなるでしょうか。
期待する性質は次のとおりです。

\[
  \texttt{rollbackAll (migrateAll xs) = xs}
\]

\[
  \texttt{migrateAll (rollbackAll ys) = ys}
\]

直感的には、各要素が戻るなら、リスト全体も戻ります。
Lean では、この直感を帰納法で証明します。
空リストは、空リストへ戻ります。
先頭要素については 1 件の round-trip 証明を使い、残りのリストについては帰納法の仮定を使います。

\begin{listing}[htbp]
\begin{minted}{lean4}
theorem rollbackAll_migrateAll (xs : List UserV1) :
    rollbackAll (migrateAll xs) = xs := by
  induction xs with
  | nil =>
      rfl
  | cons x xs ih =>
      unfold rollbackAll migrateAll at ih ⊢
      simp [rollback_migrate, ih]

theorem migrateAll_rollbackAll (ys : List UserV2) :
    migrateAll (rollbackAll ys) = ys := by
  induction ys with
  | nil =>
      rfl
  | cons y ys ih =>
      unfold rollbackAll migrateAll at ih ⊢
      simp [migrate_rollback, ih]
\end{minted}
\caption{\texttt{rollbackAll\_migrateAll} と \texttt{migrateAll\_rollbackAll}}
\label{lst:ch30-list-roundtrip}
\end{listing}

この二つが成り立つと、少なくとも Lean でモデル化した範囲では、1 件の lossless migration をリスト全体の lossless migration に拡張できることが分かります。
これは、バッチ移行の非常に重要な性質です。

図\ref{fig:ch30-roundtrip-lift}は、各要素の round-trip とリストの構造帰納法から、リスト全体の round-trip を得る構成を表します。

\FigurePlaceholder{fig-ch30-roundtrip-lift.png}{0.90}{round-trip の持ち上げ}{fig:ch30-roundtrip-lift}

\section{ID保存をリスト全体で読む}

件数保存と round-trip は重要ですが、実務レビューでは、さらに業務に近い性質も確認したくなります。
たとえば、移行前後でユーザー ID の列が保たれることを確認します。

\texttt{UserV1} では ID フィールドは \texttt{id} であり、\texttt{UserV2} では \texttt{userId} です。
個々のユーザーについて \texttt{migrate} が ID を保つなら、リスト全体でも ID 列が保たれるはずです。

\begin{listing}[htbp]
\begin{minted}{lean4}
def idsV1 (xs : List UserV1) : List Nat :=
  xs.map UserV1.id

def idsV2 (xs : List UserV2) : List Nat :=
  xs.map UserV2.userId

theorem migrateAll_preservesIds (xs : List UserV1) :
    idsV2 (migrateAll xs) = idsV1 xs := by
  induction xs with
  | nil =>
      rfl
  | cons x xs ih =>
      simp [idsV1, idsV2, migrateAll, migrate]
\end{minted}
\caption{\texttt{migrateAll\_preservesIds}}
\label{lst:ch30-id-preservation}
\end{listing}

この定理は、移行前後で ID の並びが変わらないことを表しています。
単に ID 集合が同じというより強い主張です。
順序も含めて同じ列であることを主張しているためです。
実務で順序を保証しない移行であれば、この定理は強すぎます。
その場合は、集合としての一致、重複を含む multiset としての一致、キーごとの対応など、別の仕様に落とす必要があります。

\section{バッチ移行の安全性チェックリスト}

ここまでの証明を実務へ持ち込むときは、何を \texttt{List.map} のモデルで扱い、何を別の仕組みで扱うかを分ける必要があります。
以下の表は、バッチ移行レビューで確認する項目を、Lean で扱いやすいものと、実システム側の確認が必要なものに分けています。

\begin{longtable}{p{0.28\linewidth}p{0.33\linewidth}p{0.30\linewidth}}
\toprule
確認項目 & Lean での表し方 & 実務上の注意 \\
\midrule
件数保存 & \texttt{(migrateAll xs).length = xs.length} & フィルタ、重複除去、失敗時スキップをする場合は成り立ちません。 \\
ID 保存 & \texttt{idsV2 (migrateAll xs) = idsV1 xs} & 並び順まで保証するのか、キー集合だけでよいのかを決めます。 \\
round-trip & \texttt{rollbackAll (migrateAll xs) = xs} & 完全同型のときだけ主張できます。
lossy migration では仕様を弱めます。 \\
合成保存 & \texttt{map (g . f) = map g . map f} & 純粋な値変換として切り出せる部分に限って使います。 \\
エラー処理 & \texttt{Option} や \texttt{Except} を含む型にする & 失敗時に止める、スキップする、ログに残すなどの方針を別途モデル化します。 \\
DB 更新 & 実装・運用で検査 & トランザクション、ロック、リトライ、部分失敗を統合テストや運用検証で確認します。 \\
\bottomrule
\end{longtable}

\section{本章のまとめ}

本章では、\texttt{List.map} を使って 1 件の移行をリスト全体へ持ち上げました。
確認した主張は、件数保存、合成保存、round-trip の持ち上げ、ID 列の保存です。
これらは、バッチ移行の仕様レビューでそのまま使いやすい形をしています。

失敗する移行、外部副作用、DB 更新、並び順の変更、重複除去は、\texttt{Option}、\texttt{Except}、状態、関係、集合的な仕様と、実装・運用上の検査を組み合わせて表します。
純粋なリスト変換の法則は、それらの検査を組み立てる核になります。

\texttt{migrate : UserV1 -> UserV2} は、\texttt{List.map migrate} によって \texttt{List UserV1 -> List UserV2} へ持ち上げられます。

\texttt{List.map} はリストの外側の構造を保ち、中身だけを変換します。

件数保存は、\texttt{List.map} の基本性質として帰納法で証明できます。

関手の合成保存により、二段階変換のまとめ方を変えても結果が一致します。

1 件の round-trip 性質があれば、リスト全体の round-trip 性質も帰納法で証明できます。

Lean の証明は、バッチ移行の純粋な変換核についての保証であり、実システム全体の保証とは区別します。
